% ============================================================================
% S-LoRA — cong thuc toan cua phan Method.
% Ky hieu khop fig_orientations.png: \overline{W} cho khoi dong bang.
%
% File nay TU BIEN DICH DUOC: pdflatex method.tex
%
% De ghep vao paper: xoa khoi PREAMBLE (den \begin{document}) va dong
% \end{document} cuoi file, roi % ============================================================================
% S-LoRA — cong thuc toan cua phan Method.
% Ky hieu khop fig_orientations.png: \overline{W} cho khoi dong bang.
%
% File nay TU BIEN DICH DUOC: pdflatex method.tex
%
% De ghep vao paper: xoa khoi PREAMBLE (den \begin{document}) va dong
% \end{document} cuoi file, roi % ============================================================================
% S-LoRA — cong thuc toan cua phan Method.
% Ky hieu khop fig_orientations.png: \overline{W} cho khoi dong bang.
%
% File nay TU BIEN DICH DUOC: pdflatex method.tex
%
% De ghep vao paper: xoa khoi PREAMBLE (den \begin{document}) va dong
% \end{document} cuoi file, roi % ============================================================================
% S-LoRA — cong thuc toan cua phan Method.
% Ky hieu khop fig_orientations.png: \overline{W} cho khoi dong bang.
%
% File nay TU BIEN DICH DUOC: pdflatex method.tex
%
% De ghep vao paper: xoa khoi PREAMBLE (den \begin{document}) va dong
% \end{document} cuoi file, roi \input{method} tu file chinh. Cac \providecommand
% giu nguyen duoc — chung khong ghi de dinh nghia da co.
% Nhan dat tien to "eq:m-" nen khong dung voi nhan trong iclr_draft.tex.
% ============================================================================
\documentclass[11pt]{article}
\usepackage{amsmath,amssymb}
\usepackage[margin=1in]{geometry}

\providecommand{\Wb}{\overline{W}}
\providecommand{\R}{\mathbb{R}}
\providecommand{\rowsp}{\operatorname{row}}
\providecommand{\colsp}{\operatorname{col}}

\begin{document}
% ==================== PREAMBLE KET THUC — phan duoi la noi dung ============

\section*{S-LoRA: the factorization}

% ---------------------------------------------------------------- ky hieu
Let $W_0\in\R^{n\times m}$ with $n\neq m$, acting as $y=W_0x$, and set
\begin{equation}
  k=\min(n,m),\qquad a=\frac{\max(n,m)}{\min(n,m)}\ge 1,
  \qquad nm=ak^2 .
  \label{eq:m-notation}
\end{equation}

% ---------------------------------------------------------------- beo ngang
\paragraph{Wide weights ($m>n$, $k=n$): split the columns.}
With $\Pi$ a column permutation,
\begin{equation}
  W_0\Pi=\begin{bmatrix}W_1 & W_2\end{bmatrix},\qquad
  W_1\in\R^{n\times(m-n)},\quad W_2\in\R^{n\times n},
  \label{eq:m-split-wide}
\end{equation}
\begin{equation}
  W_2\Wb=W_1
  \quad\Longrightarrow\quad
  \Wb=W_2^{-1}W_1\in\R^{n\times(m-n)},
  \label{eq:m-solve-wide}
\end{equation}
\begin{equation}
  W_0\Pi=\begin{bmatrix}W_2\Wb & W_2\end{bmatrix}
        =W_2\begin{bmatrix}\Wb & I\end{bmatrix}.
  \label{eq:m-wide}
\end{equation}
For $\Pi^{\!\top}x=[\,x_1;x_2\,]$ with $x_1\in\R^{m-n}$, $x_2\in\R^{n}$,
\begin{equation}
  W_0x=W_1x_1+W_2x_2=W_2\!\left(\Wb x_1+x_2\right).
  \label{eq:m-fwd-wide}
\end{equation}

% ---------------------------------------------------------------- cao doc
\paragraph{Tall weights ($m<n$, $k=m$): split the rows.}
With $\Pi$ a row permutation,
\begin{equation}
  \Pi W_0=\begin{bmatrix}W_1\\ W_2\end{bmatrix},\qquad
  W_1\in\R^{(n-m)\times m},\quad W_2\in\R^{m\times m},
  \label{eq:m-split-tall}
\end{equation}
\begin{equation}
  \Wb W_2=W_1
  \quad\Longrightarrow\quad
  \Wb=W_1W_2^{-1}\in\R^{(n-m)\times m},
  \label{eq:m-solve-tall}
\end{equation}
\begin{equation}
  \Pi W_0=\begin{bmatrix}\Wb W_2\\ W_2\end{bmatrix}
        =\begin{bmatrix}\Wb\\ I\end{bmatrix}W_2,
  \qquad
  \Pi W_0x=\begin{bmatrix}W_1x\\ W_2x\end{bmatrix}
          =\begin{bmatrix}\Wb\\ I\end{bmatrix}W_2x .
  \label{eq:m-tall}
\end{equation}

% ---------------------------------------------------------------- rang buoc
\paragraph{The constraint.}
Write $P=[\,\Wb\mid I\,]\Pi^{\!\top}$ for $m>n$ and
$P=\Pi^{\!\top}[\,\Wb\,;I\,]$ for $m<n$, so that $W_0=W_2P$ or $W_0=PW_2$.
Freezing $\Wb$ and perturbing $W_2\mapsto W_2+\Delta W_2$ gives
\begin{align}
  m>n:\quad \Delta W&=\Delta W_2\,P
    && \Longrightarrow\ \rowsp(\Delta W)\subseteq\rowsp(W_0),
    \label{eq:m-row}\\
  m<n:\quad \Delta W&=P\,\Delta W_2
    && \Longrightarrow\ \colsp(\Delta W)\subseteq\colsp(W_0),
    \label{eq:m-col}
\end{align}
since $W_2$ nonsingular gives $\rowsp(P)=\rowsp(W_2P)=\rowsp(W_0)$, and dually
for columns. The containment holds for \emph{every} $\Delta W_2$.

% ---------------------------------------------------------------- pivoted QR
\paragraph{Choosing $\Pi$: column-pivoted QR.}
\begin{equation}
  W_0\Pi=Q\begin{bmatrix}R_{11} & R_{12}\end{bmatrix},\qquad
  Q^{\!\top}Q=I,\quad R_{11}\in\R^{n\times n}\ \text{upper triangular},
  \label{eq:m-qr}
\end{equation}
so that $W_2=QR_{11}$, $W_1=QR_{12}$, and the orthogonal factor cancels:
\begin{equation}
  \Wb=W_2^{-1}W_1=(QR_{11})^{-1}(QR_{12})=R_{11}^{-1}R_{12}.
  \label{eq:m-qcancel}
\end{equation}
One triangular solve; $Q$ is never formed and no inverse is computed.

% ---------------------------------------------------------------- ke toan
\paragraph{Parameter and cost accounting.}
Storing $\Wb$ and $W_2$ (the identity block is an index selection) costs
\begin{equation}
  \underbrace{k\left|m-n\right|}_{\Wb}+\underbrace{k^{2}}_{W_2}
  =k\!\left(\left|m-n\right|+k\right)=nm ,
  \label{eq:m-preserve}
\end{equation}
so the parameter count is preserved exactly. A rank-$r$ adapter whose cost
scales with the dimensions of its host matrix costs $r(n+m)=rk(1+a)$ on $W_0$
but $2rk$ on $W_2$, hence
\begin{equation}
  \frac{r(n+m)}{2rk}=\frac{1+a}{2},
  \qquad\text{and for full fine-tuning}\qquad
  \frac{nm}{k^{2}}=a .
  \label{eq:m-advantage}
\end{equation}

% ---------------------------------------------------------------- LoRA han che
\paragraph{S-LoRA as restricted LoRA.}
With $\Delta W_2=\tfrac{\alpha}{r}BA$, $B\in\R^{k\times r}$,
$A\in\R^{r\times k}$ and $B$ zero-initialized, \eqref{eq:m-wide} gives
\begin{equation}
  \Delta W=\tfrac{\alpha}{r}\,B\,(AP),
  \qquad AP\in\R^{r\times m},
  \label{eq:m-restricted}
\end{equation}
a rank-$r$ update whose down-projection is trainable but confined to
$\rowsp(W_0)$.

% ---------------------------------------------------------------- gradient
\paragraph{Gradient view.}
Holding $P$ fixed, $\partial L/\partial W_2=(\partial L/\partial W)P^{\!\top}$,
so a step $\Delta W_2=-\eta\,(\partial L/\partial W)P^{\!\top}$ induces
\begin{equation}
  \Delta W=\Delta W_2P=-\eta\,\frac{\partial L}{\partial W}\,P^{\!\top}\!P .
  \label{eq:m-precond}
\end{equation}
If $P$ had orthonormal rows, $P^{\!\top}\!P=VV^{\!\top}$ would be the orthogonal
projector onto $\rowsp(W_0)$; with the pivoted-QR basis it projects and
rescales.

% ---------------------------------------------------------------- merge
\paragraph{Merging.}
Let $W_2'=W_2+\Delta W_2$ be the trained square factor. Then
\begin{equation}
  W=W_2'\begin{bmatrix}\Wb & I\end{bmatrix}\Pi^{\!\top}
  \qquad\text{or}\qquad
  W=\Pi^{\!\top}\begin{bmatrix}\Wb\\ I\end{bmatrix}W_2'
  \label{eq:m-merge}
\end{equation}
recovers a single matrix of the original shape.

% ==================== XOA TU DAY XUONG khi ghep vao paper =================
\end{document}
 tu file chinh. Cac \providecommand
% giu nguyen duoc — chung khong ghi de dinh nghia da co.
% Nhan dat tien to "eq:m-" nen khong dung voi nhan trong iclr_draft.tex.
% ============================================================================
\documentclass[11pt]{article}
\usepackage{amsmath,amssymb}
\usepackage[margin=1in]{geometry}

\providecommand{\Wb}{\overline{W}}
\providecommand{\R}{\mathbb{R}}
\providecommand{\rowsp}{\operatorname{row}}
\providecommand{\colsp}{\operatorname{col}}

\begin{document}
% ==================== PREAMBLE KET THUC — phan duoi la noi dung ============

\section*{S-LoRA: the factorization}

% ---------------------------------------------------------------- ky hieu
Let $W_0\in\R^{n\times m}$ with $n\neq m$, acting as $y=W_0x$, and set
\begin{equation}
  k=\min(n,m),\qquad a=\frac{\max(n,m)}{\min(n,m)}\ge 1,
  \qquad nm=ak^2 .
  \label{eq:m-notation}
\end{equation}

% ---------------------------------------------------------------- beo ngang
\paragraph{Wide weights ($m>n$, $k=n$): split the columns.}
With $\Pi$ a column permutation,
\begin{equation}
  W_0\Pi=\begin{bmatrix}W_1 & W_2\end{bmatrix},\qquad
  W_1\in\R^{n\times(m-n)},\quad W_2\in\R^{n\times n},
  \label{eq:m-split-wide}
\end{equation}
\begin{equation}
  W_2\Wb=W_1
  \quad\Longrightarrow\quad
  \Wb=W_2^{-1}W_1\in\R^{n\times(m-n)},
  \label{eq:m-solve-wide}
\end{equation}
\begin{equation}
  W_0\Pi=\begin{bmatrix}W_2\Wb & W_2\end{bmatrix}
        =W_2\begin{bmatrix}\Wb & I\end{bmatrix}.
  \label{eq:m-wide}
\end{equation}
For $\Pi^{\!\top}x=[\,x_1;x_2\,]$ with $x_1\in\R^{m-n}$, $x_2\in\R^{n}$,
\begin{equation}
  W_0x=W_1x_1+W_2x_2=W_2\!\left(\Wb x_1+x_2\right).
  \label{eq:m-fwd-wide}
\end{equation}

% ---------------------------------------------------------------- cao doc
\paragraph{Tall weights ($m<n$, $k=m$): split the rows.}
With $\Pi$ a row permutation,
\begin{equation}
  \Pi W_0=\begin{bmatrix}W_1\\ W_2\end{bmatrix},\qquad
  W_1\in\R^{(n-m)\times m},\quad W_2\in\R^{m\times m},
  \label{eq:m-split-tall}
\end{equation}
\begin{equation}
  \Wb W_2=W_1
  \quad\Longrightarrow\quad
  \Wb=W_1W_2^{-1}\in\R^{(n-m)\times m},
  \label{eq:m-solve-tall}
\end{equation}
\begin{equation}
  \Pi W_0=\begin{bmatrix}\Wb W_2\\ W_2\end{bmatrix}
        =\begin{bmatrix}\Wb\\ I\end{bmatrix}W_2,
  \qquad
  \Pi W_0x=\begin{bmatrix}W_1x\\ W_2x\end{bmatrix}
          =\begin{bmatrix}\Wb\\ I\end{bmatrix}W_2x .
  \label{eq:m-tall}
\end{equation}

% ---------------------------------------------------------------- rang buoc
\paragraph{The constraint.}
Write $P=[\,\Wb\mid I\,]\Pi^{\!\top}$ for $m>n$ and
$P=\Pi^{\!\top}[\,\Wb\,;I\,]$ for $m<n$, so that $W_0=W_2P$ or $W_0=PW_2$.
Freezing $\Wb$ and perturbing $W_2\mapsto W_2+\Delta W_2$ gives
\begin{align}
  m>n:\quad \Delta W&=\Delta W_2\,P
    && \Longrightarrow\ \rowsp(\Delta W)\subseteq\rowsp(W_0),
    \label{eq:m-row}\\
  m<n:\quad \Delta W&=P\,\Delta W_2
    && \Longrightarrow\ \colsp(\Delta W)\subseteq\colsp(W_0),
    \label{eq:m-col}
\end{align}
since $W_2$ nonsingular gives $\rowsp(P)=\rowsp(W_2P)=\rowsp(W_0)$, and dually
for columns. The containment holds for \emph{every} $\Delta W_2$.

% ---------------------------------------------------------------- pivoted QR
\paragraph{Choosing $\Pi$: column-pivoted QR.}
\begin{equation}
  W_0\Pi=Q\begin{bmatrix}R_{11} & R_{12}\end{bmatrix},\qquad
  Q^{\!\top}Q=I,\quad R_{11}\in\R^{n\times n}\ \text{upper triangular},
  \label{eq:m-qr}
\end{equation}
so that $W_2=QR_{11}$, $W_1=QR_{12}$, and the orthogonal factor cancels:
\begin{equation}
  \Wb=W_2^{-1}W_1=(QR_{11})^{-1}(QR_{12})=R_{11}^{-1}R_{12}.
  \label{eq:m-qcancel}
\end{equation}
One triangular solve; $Q$ is never formed and no inverse is computed.

% ---------------------------------------------------------------- ke toan
\paragraph{Parameter and cost accounting.}
Storing $\Wb$ and $W_2$ (the identity block is an index selection) costs
\begin{equation}
  \underbrace{k\left|m-n\right|}_{\Wb}+\underbrace{k^{2}}_{W_2}
  =k\!\left(\left|m-n\right|+k\right)=nm ,
  \label{eq:m-preserve}
\end{equation}
so the parameter count is preserved exactly. A rank-$r$ adapter whose cost
scales with the dimensions of its host matrix costs $r(n+m)=rk(1+a)$ on $W_0$
but $2rk$ on $W_2$, hence
\begin{equation}
  \frac{r(n+m)}{2rk}=\frac{1+a}{2},
  \qquad\text{and for full fine-tuning}\qquad
  \frac{nm}{k^{2}}=a .
  \label{eq:m-advantage}
\end{equation}

% ---------------------------------------------------------------- LoRA han che
\paragraph{S-LoRA as restricted LoRA.}
With $\Delta W_2=\tfrac{\alpha}{r}BA$, $B\in\R^{k\times r}$,
$A\in\R^{r\times k}$ and $B$ zero-initialized, \eqref{eq:m-wide} gives
\begin{equation}
  \Delta W=\tfrac{\alpha}{r}\,B\,(AP),
  \qquad AP\in\R^{r\times m},
  \label{eq:m-restricted}
\end{equation}
a rank-$r$ update whose down-projection is trainable but confined to
$\rowsp(W_0)$.

% ---------------------------------------------------------------- gradient
\paragraph{Gradient view.}
Holding $P$ fixed, $\partial L/\partial W_2=(\partial L/\partial W)P^{\!\top}$,
so a step $\Delta W_2=-\eta\,(\partial L/\partial W)P^{\!\top}$ induces
\begin{equation}
  \Delta W=\Delta W_2P=-\eta\,\frac{\partial L}{\partial W}\,P^{\!\top}\!P .
  \label{eq:m-precond}
\end{equation}
If $P$ had orthonormal rows, $P^{\!\top}\!P=VV^{\!\top}$ would be the orthogonal
projector onto $\rowsp(W_0)$; with the pivoted-QR basis it projects and
rescales.

% ---------------------------------------------------------------- merge
\paragraph{Merging.}
Let $W_2'=W_2+\Delta W_2$ be the trained square factor. Then
\begin{equation}
  W=W_2'\begin{bmatrix}\Wb & I\end{bmatrix}\Pi^{\!\top}
  \qquad\text{or}\qquad
  W=\Pi^{\!\top}\begin{bmatrix}\Wb\\ I\end{bmatrix}W_2'
  \label{eq:m-merge}
\end{equation}
recovers a single matrix of the original shape.

% ==================== XOA TU DAY XUONG khi ghep vao paper =================
\end{document}
 tu file chinh. Cac \providecommand
% giu nguyen duoc — chung khong ghi de dinh nghia da co.
% Nhan dat tien to "eq:m-" nen khong dung voi nhan trong iclr_draft.tex.
% ============================================================================
\documentclass[11pt]{article}
\usepackage{amsmath,amssymb}
\usepackage[margin=1in]{geometry}

\providecommand{\Wb}{\overline{W}}
\providecommand{\R}{\mathbb{R}}
\providecommand{\rowsp}{\operatorname{row}}
\providecommand{\colsp}{\operatorname{col}}

\begin{document}
% ==================== PREAMBLE KET THUC — phan duoi la noi dung ============

\section*{S-LoRA: the factorization}

% ---------------------------------------------------------------- ky hieu
Let $W_0\in\R^{n\times m}$ with $n\neq m$, acting as $y=W_0x$, and set
\begin{equation}
  k=\min(n,m),\qquad a=\frac{\max(n,m)}{\min(n,m)}\ge 1,
  \qquad nm=ak^2 .
  \label{eq:m-notation}
\end{equation}

% ---------------------------------------------------------------- beo ngang
\paragraph{Wide weights ($m>n$, $k=n$): split the columns.}
With $\Pi$ a column permutation,
\begin{equation}
  W_0\Pi=\begin{bmatrix}W_1 & W_2\end{bmatrix},\qquad
  W_1\in\R^{n\times(m-n)},\quad W_2\in\R^{n\times n},
  \label{eq:m-split-wide}
\end{equation}
\begin{equation}
  W_2\Wb=W_1
  \quad\Longrightarrow\quad
  \Wb=W_2^{-1}W_1\in\R^{n\times(m-n)},
  \label{eq:m-solve-wide}
\end{equation}
\begin{equation}
  W_0\Pi=\begin{bmatrix}W_2\Wb & W_2\end{bmatrix}
        =W_2\begin{bmatrix}\Wb & I\end{bmatrix}.
  \label{eq:m-wide}
\end{equation}
For $\Pi^{\!\top}x=[\,x_1;x_2\,]$ with $x_1\in\R^{m-n}$, $x_2\in\R^{n}$,
\begin{equation}
  W_0x=W_1x_1+W_2x_2=W_2\!\left(\Wb x_1+x_2\right).
  \label{eq:m-fwd-wide}
\end{equation}

% ---------------------------------------------------------------- cao doc
\paragraph{Tall weights ($m<n$, $k=m$): split the rows.}
With $\Pi$ a row permutation,
\begin{equation}
  \Pi W_0=\begin{bmatrix}W_1\\ W_2\end{bmatrix},\qquad
  W_1\in\R^{(n-m)\times m},\quad W_2\in\R^{m\times m},
  \label{eq:m-split-tall}
\end{equation}
\begin{equation}
  \Wb W_2=W_1
  \quad\Longrightarrow\quad
  \Wb=W_1W_2^{-1}\in\R^{(n-m)\times m},
  \label{eq:m-solve-tall}
\end{equation}
\begin{equation}
  \Pi W_0=\begin{bmatrix}\Wb W_2\\ W_2\end{bmatrix}
        =\begin{bmatrix}\Wb\\ I\end{bmatrix}W_2,
  \qquad
  \Pi W_0x=\begin{bmatrix}W_1x\\ W_2x\end{bmatrix}
          =\begin{bmatrix}\Wb\\ I\end{bmatrix}W_2x .
  \label{eq:m-tall}
\end{equation}

% ---------------------------------------------------------------- rang buoc
\paragraph{The constraint.}
Write $P=[\,\Wb\mid I\,]\Pi^{\!\top}$ for $m>n$ and
$P=\Pi^{\!\top}[\,\Wb\,;I\,]$ for $m<n$, so that $W_0=W_2P$ or $W_0=PW_2$.
Freezing $\Wb$ and perturbing $W_2\mapsto W_2+\Delta W_2$ gives
\begin{align}
  m>n:\quad \Delta W&=\Delta W_2\,P
    && \Longrightarrow\ \rowsp(\Delta W)\subseteq\rowsp(W_0),
    \label{eq:m-row}\\
  m<n:\quad \Delta W&=P\,\Delta W_2
    && \Longrightarrow\ \colsp(\Delta W)\subseteq\colsp(W_0),
    \label{eq:m-col}
\end{align}
since $W_2$ nonsingular gives $\rowsp(P)=\rowsp(W_2P)=\rowsp(W_0)$, and dually
for columns. The containment holds for \emph{every} $\Delta W_2$.

% ---------------------------------------------------------------- pivoted QR
\paragraph{Choosing $\Pi$: column-pivoted QR.}
\begin{equation}
  W_0\Pi=Q\begin{bmatrix}R_{11} & R_{12}\end{bmatrix},\qquad
  Q^{\!\top}Q=I,\quad R_{11}\in\R^{n\times n}\ \text{upper triangular},
  \label{eq:m-qr}
\end{equation}
so that $W_2=QR_{11}$, $W_1=QR_{12}$, and the orthogonal factor cancels:
\begin{equation}
  \Wb=W_2^{-1}W_1=(QR_{11})^{-1}(QR_{12})=R_{11}^{-1}R_{12}.
  \label{eq:m-qcancel}
\end{equation}
One triangular solve; $Q$ is never formed and no inverse is computed.

% ---------------------------------------------------------------- ke toan
\paragraph{Parameter and cost accounting.}
Storing $\Wb$ and $W_2$ (the identity block is an index selection) costs
\begin{equation}
  \underbrace{k\left|m-n\right|}_{\Wb}+\underbrace{k^{2}}_{W_2}
  =k\!\left(\left|m-n\right|+k\right)=nm ,
  \label{eq:m-preserve}
\end{equation}
so the parameter count is preserved exactly. A rank-$r$ adapter whose cost
scales with the dimensions of its host matrix costs $r(n+m)=rk(1+a)$ on $W_0$
but $2rk$ on $W_2$, hence
\begin{equation}
  \frac{r(n+m)}{2rk}=\frac{1+a}{2},
  \qquad\text{and for full fine-tuning}\qquad
  \frac{nm}{k^{2}}=a .
  \label{eq:m-advantage}
\end{equation}

% ---------------------------------------------------------------- LoRA han che
\paragraph{S-LoRA as restricted LoRA.}
With $\Delta W_2=\tfrac{\alpha}{r}BA$, $B\in\R^{k\times r}$,
$A\in\R^{r\times k}$ and $B$ zero-initialized, \eqref{eq:m-wide} gives
\begin{equation}
  \Delta W=\tfrac{\alpha}{r}\,B\,(AP),
  \qquad AP\in\R^{r\times m},
  \label{eq:m-restricted}
\end{equation}
a rank-$r$ update whose down-projection is trainable but confined to
$\rowsp(W_0)$.

% ---------------------------------------------------------------- gradient
\paragraph{Gradient view.}
Holding $P$ fixed, $\partial L/\partial W_2=(\partial L/\partial W)P^{\!\top}$,
so a step $\Delta W_2=-\eta\,(\partial L/\partial W)P^{\!\top}$ induces
\begin{equation}
  \Delta W=\Delta W_2P=-\eta\,\frac{\partial L}{\partial W}\,P^{\!\top}\!P .
  \label{eq:m-precond}
\end{equation}
If $P$ had orthonormal rows, $P^{\!\top}\!P=VV^{\!\top}$ would be the orthogonal
projector onto $\rowsp(W_0)$; with the pivoted-QR basis it projects and
rescales.

% ---------------------------------------------------------------- merge
\paragraph{Merging.}
Let $W_2'=W_2+\Delta W_2$ be the trained square factor. Then
\begin{equation}
  W=W_2'\begin{bmatrix}\Wb & I\end{bmatrix}\Pi^{\!\top}
  \qquad\text{or}\qquad
  W=\Pi^{\!\top}\begin{bmatrix}\Wb\\ I\end{bmatrix}W_2'
  \label{eq:m-merge}
\end{equation}
recovers a single matrix of the original shape.

% ==================== XOA TU DAY XUONG khi ghep vao paper =================
\end{document}
 tu file chinh. Cac \providecommand
% giu nguyen duoc — chung khong ghi de dinh nghia da co.
% Nhan dat tien to "eq:m-" nen khong dung voi nhan trong iclr_draft.tex.
% ============================================================================
\documentclass[11pt]{article}
\usepackage{amsmath,amssymb}
\usepackage[margin=1in]{geometry}

\providecommand{\Wb}{\overline{W}}
\providecommand{\R}{\mathbb{R}}
\providecommand{\rowsp}{\operatorname{row}}
\providecommand{\colsp}{\operatorname{col}}

\begin{document}
% ==================== PREAMBLE KET THUC — phan duoi la noi dung ============

\section*{S-LoRA: the factorization}

% ---------------------------------------------------------------- ky hieu
Let $W_0\in\R^{n\times m}$ with $n\neq m$, acting as $y=W_0x$, and set
\begin{equation}
  k=\min(n,m),\qquad a=\frac{\max(n,m)}{\min(n,m)}\ge 1,
  \qquad nm=ak^2 .
  \label{eq:m-notation}
\end{equation}

% ---------------------------------------------------------------- beo ngang
\paragraph{Wide weights ($m>n$, $k=n$): split the columns.}
With $\Pi$ a column permutation,
\begin{equation}
  W_0\Pi=\begin{bmatrix}W_1 & W_2\end{bmatrix},\qquad
  W_1\in\R^{n\times(m-n)},\quad W_2\in\R^{n\times n},
  \label{eq:m-split-wide}
\end{equation}
\begin{equation}
  W_2\Wb=W_1
  \quad\Longrightarrow\quad
  \Wb=W_2^{-1}W_1\in\R^{n\times(m-n)},
  \label{eq:m-solve-wide}
\end{equation}
\begin{equation}
  W_0\Pi=\begin{bmatrix}W_2\Wb & W_2\end{bmatrix}
        =W_2\begin{bmatrix}\Wb & I\end{bmatrix}.
  \label{eq:m-wide}
\end{equation}
For $\Pi^{\!\top}x=[\,x_1;x_2\,]$ with $x_1\in\R^{m-n}$, $x_2\in\R^{n}$,
\begin{equation}
  W_0x=W_1x_1+W_2x_2=W_2\!\left(\Wb x_1+x_2\right).
  \label{eq:m-fwd-wide}
\end{equation}

% ---------------------------------------------------------------- cao doc
\paragraph{Tall weights ($m<n$, $k=m$): split the rows.}
With $\Pi$ a row permutation,
\begin{equation}
  \Pi W_0=\begin{bmatrix}W_1\\ W_2\end{bmatrix},\qquad
  W_1\in\R^{(n-m)\times m},\quad W_2\in\R^{m\times m},
  \label{eq:m-split-tall}
\end{equation}
\begin{equation}
  \Wb W_2=W_1
  \quad\Longrightarrow\quad
  \Wb=W_1W_2^{-1}\in\R^{(n-m)\times m},
  \label{eq:m-solve-tall}
\end{equation}
\begin{equation}
  \Pi W_0=\begin{bmatrix}\Wb W_2\\ W_2\end{bmatrix}
        =\begin{bmatrix}\Wb\\ I\end{bmatrix}W_2,
  \qquad
  \Pi W_0x=\begin{bmatrix}W_1x\\ W_2x\end{bmatrix}
          =\begin{bmatrix}\Wb\\ I\end{bmatrix}W_2x .
  \label{eq:m-tall}
\end{equation}

% ---------------------------------------------------------------- rang buoc
\paragraph{The constraint.}
Write $P=[\,\Wb\mid I\,]\Pi^{\!\top}$ for $m>n$ and
$P=\Pi^{\!\top}[\,\Wb\,;I\,]$ for $m<n$, so that $W_0=W_2P$ or $W_0=PW_2$.
Freezing $\Wb$ and perturbing $W_2\mapsto W_2+\Delta W_2$ gives
\begin{align}
  m>n:\quad \Delta W&=\Delta W_2\,P
    && \Longrightarrow\ \rowsp(\Delta W)\subseteq\rowsp(W_0),
    \label{eq:m-row}\\
  m<n:\quad \Delta W&=P\,\Delta W_2
    && \Longrightarrow\ \colsp(\Delta W)\subseteq\colsp(W_0),
    \label{eq:m-col}
\end{align}
since $W_2$ nonsingular gives $\rowsp(P)=\rowsp(W_2P)=\rowsp(W_0)$, and dually
for columns. The containment holds for \emph{every} $\Delta W_2$.

% ---------------------------------------------------------------- pivoted QR
\paragraph{Choosing $\Pi$: column-pivoted QR.}
\begin{equation}
  W_0\Pi=Q\begin{bmatrix}R_{11} & R_{12}\end{bmatrix},\qquad
  Q^{\!\top}Q=I,\quad R_{11}\in\R^{n\times n}\ \text{upper triangular},
  \label{eq:m-qr}
\end{equation}
so that $W_2=QR_{11}$, $W_1=QR_{12}$, and the orthogonal factor cancels:
\begin{equation}
  \Wb=W_2^{-1}W_1=(QR_{11})^{-1}(QR_{12})=R_{11}^{-1}R_{12}.
  \label{eq:m-qcancel}
\end{equation}
One triangular solve; $Q$ is never formed and no inverse is computed.

% ---------------------------------------------------------------- ke toan
\paragraph{Parameter and cost accounting.}
Storing $\Wb$ and $W_2$ (the identity block is an index selection) costs
\begin{equation}
  \underbrace{k\left|m-n\right|}_{\Wb}+\underbrace{k^{2}}_{W_2}
  =k\!\left(\left|m-n\right|+k\right)=nm ,
  \label{eq:m-preserve}
\end{equation}
so the parameter count is preserved exactly. A rank-$r$ adapter whose cost
scales with the dimensions of its host matrix costs $r(n+m)=rk(1+a)$ on $W_0$
but $2rk$ on $W_2$, hence
\begin{equation}
  \frac{r(n+m)}{2rk}=\frac{1+a}{2},
  \qquad\text{and for full fine-tuning}\qquad
  \frac{nm}{k^{2}}=a .
  \label{eq:m-advantage}
\end{equation}

% ---------------------------------------------------------------- LoRA han che
\paragraph{S-LoRA as restricted LoRA.}
With $\Delta W_2=\tfrac{\alpha}{r}BA$, $B\in\R^{k\times r}$,
$A\in\R^{r\times k}$ and $B$ zero-initialized, \eqref{eq:m-wide} gives
\begin{equation}
  \Delta W=\tfrac{\alpha}{r}\,B\,(AP),
  \qquad AP\in\R^{r\times m},
  \label{eq:m-restricted}
\end{equation}
a rank-$r$ update whose down-projection is trainable but confined to
$\rowsp(W_0)$.

% ---------------------------------------------------------------- gradient
\paragraph{Gradient view.}
Holding $P$ fixed, $\partial L/\partial W_2=(\partial L/\partial W)P^{\!\top}$,
so a step $\Delta W_2=-\eta\,(\partial L/\partial W)P^{\!\top}$ induces
\begin{equation}
  \Delta W=\Delta W_2P=-\eta\,\frac{\partial L}{\partial W}\,P^{\!\top}\!P .
  \label{eq:m-precond}
\end{equation}
If $P$ had orthonormal rows, $P^{\!\top}\!P=VV^{\!\top}$ would be the orthogonal
projector onto $\rowsp(W_0)$; with the pivoted-QR basis it projects and
rescales.

% ---------------------------------------------------------------- merge
\paragraph{Merging.}
Let $W_2'=W_2+\Delta W_2$ be the trained square factor. Then
\begin{equation}
  W=W_2'\begin{bmatrix}\Wb & I\end{bmatrix}\Pi^{\!\top}
  \qquad\text{or}\qquad
  W=\Pi^{\!\top}\begin{bmatrix}\Wb\\ I\end{bmatrix}W_2'
  \label{eq:m-merge}
\end{equation}
recovers a single matrix of the original shape.

% ==================== XOA TU DAY XUONG khi ghep vao paper =================
\end{document}
